\section{Method}
\label{sec:method}

\subsection{FT-Transformer}
\label{sec:ft_transformer}

We adopt the FT-Transformer architecture of \citet{gorishniy2021revisiting} as our primary encoder for structured financial data. The architecture consists of three components:

\paragraph{Feature Tokenizer.} Each scalar input feature $x_i$ is projected to a $d$-dimensional token embedding via a per-feature affine transformation: $\mathbf{t}_i = \mathbf{W}_i x_i + \mathbf{b}_i$, where $\mathbf{W}_i \in \mathbb{R}^d$ and $\mathbf{b}_i \in \mathbb{R}^d$ are learnable parameters. Feature weights are initialized with Kaiming uniform initialization ($a = \sqrt{5}$). A learnable \textsc{[CLS]} token $\mathbf{t}_{\mathrm{cls}} \in \mathbb{R}^d$ is prepended to the token sequence, yielding a sequence of $n + 1$ tokens where $n$ is the number of input features.

\paragraph{Transformer Encoder.} The token sequence is processed by a stack of $L$ standard Transformer encoder layers \citep{vaswani2017attention}. Each layer applies multi-head self-attention followed by a position-wise feed-forward network, with pre-layer normalization (norm-first configuration). The feed-forward hidden dimension is set to $d \times f$ where $f$ is a configurable expansion factor. We use PyTorch's \texttt{nn.TransformerEncoderLayer} with batch-first processing.

\paragraph{Classification Head.} The \textsc{[CLS]} token representation from the final encoder layer is passed through a classification head consisting of LayerNorm $\rightarrow$ ReLU $\rightarrow$ Linear, producing a single logit for binary classification. The model is trained with binary cross-entropy with logits loss (\texttt{BCEWithLogitsLoss}).

Default hyperparameters: $d = 192$ (token dimension), $h = 8$ (attention heads), $L = 3$ (layers), $f = 4$ (FFN expansion factor), dropout $= 0.0$. These follow the recommendations in \citet{gorishniy2021revisiting} for medium-sized tabular datasets.

\paragraph{Hyperparameter selection.} We conduct a grid search over 36 configurations: learning rate $\in \{10^{-3}, 3{\times}10^{-4}, 10^{-4}, 3{\times}10^{-5}\}$, number of layers $L \in \{2, 3, 4\}$, and token dimension $d \in \{128, 192, 256\}$, with the remaining hyperparameters held at their defaults. The best configuration is selected by mean test AUROC across all evaluable outcomes. All other training settings (optimizer, early stopping, batch size) are held constant across the sweep.

\subsection{Baseline Models}
\label{sec:baselines}

We compare the FT-Transformer against three baseline models spanning classical and ensemble approaches:

\paragraph{Logistic Regression (LR).} An L2-regularized logistic regression with balanced class weighting, preceded by standard scaling (scikit-learn \texttt{StandardScaler} $\rightarrow$ \texttt{LogisticRegression}). Regularization strength $C = 1.0$, solver = L-BFGS, maximum 1{,}000 iterations. We evaluate two feature configurations: LR-full (all engineered features) and LR-traditional (10 traditional financial ratios only: Altman Z-score, current ratio, quick ratio, debt-to-equity, debt-to-assets, ROA, ROE, net profit margin, interest coverage, and CFO-to-debt).

\paragraph{XGBoost.} Extreme gradient-boosted trees \citep{chen2016xgboost} with 500 estimators, learning rate 0.05, max depth~6, subsample ratio 0.8, column subsample ratio 0.8. Evaluation metric: AUC.

\paragraph{Histogram Gradient-Boosted Trees (GBT).} Scikit-learn's \texttt{HistGradientBoostingClassifier} with 500 iterations, learning rate 0.05, max depth~6, minimum 20 samples per leaf, balanced class weighting. This model natively handles NaN values, making it suitable for evaluation on raw XBRL features with minimal preprocessing.

All baseline hyperparameters use default (untuned) values for the initial benchmark. Optional Optuna-based temporal cross-validation tuning is supported (30 trials, 3-fold expanding-window CV) but is disabled in the default configuration.

\subsection{Self-Supervised Pretraining}
\label{sec:ssl}

We investigate masked feature reconstruction as a self-supervised pretraining objective for the FT-Transformer, analogous to masked language modeling \citep{devlin2019bert} applied to tabular features.

\paragraph{Masking strategy.} For each training sample, a fraction $r$ of input features are independently selected for masking via Bernoulli sampling. Masked feature values are replaced with zero in the input. We experiment with mask ratios $r \in \{0.15, 0.30, 0.50\}$.

\paragraph{Reconstruction head.} A two-layer MLP reconstructs all $n$ feature values from the \textsc{[CLS]} token representation: $\text{Linear}(d, 2d) \rightarrow \text{GELU} \rightarrow \text{Linear}(2d, n)$. The reconstruction loss is computed as the mean squared error over masked positions only:
\begin{equation}
  \mathcal{L}_{\text{SSL}} = \frac{\sum_{i \in M} (\hat{x}_i - x_i)^2}{|M|}
  \label{eq:ssl_loss}
\end{equation}
where $M$ is the set of masked feature indices and $\hat{x}_i$ is the predicted value for feature~$i$.

\paragraph{Pretraining protocol.} SSL pretraining uses the best FT-Transformer architecture identified by the hyperparameter sweep (\Cref{sec:ft_transformer}). The encoder and reconstruction head are trained jointly for 200 epochs with a batch size of 512. We use AdamW optimization with learning rate $10^{-4}$ and weight decay $10^{-5}$. The learning rate follows a linear warmup over 10 epochs followed by cosine annealing to zero. Pretraining uses all available training-split data regardless of label availability.

\paragraph{Fine-tuning.} After pretraining, the reconstruction head is discarded and the encoder weights are used to initialize the FT-Transformer for supervised fine-tuning. Fine-tuning uses the same protocol as training from scratch (\Cref{sec:training}).

\subsection{Training Protocol}
\label{sec:training}

\paragraph{FT-Transformer training.} Batch size 256, maximum 100 epochs, AdamW optimizer with learning rate $10^{-4}$ and weight decay $10^{-5}$. Early stopping with patience 10 epochs on validation AUROC. Gradient norms are clipped to 1.0. Predictions are obtained by applying sigmoid to model logits.

\paragraph{Baseline training.} Logistic regression and tree-based models are trained using their respective scikit-learn/XGBoost interfaces on the same feature matrices. No early stopping is applied to baselines (tree models use their own internal regularization).

\paragraph{Reproducibility.} All experiments use random seed 42 for weight initialization, data shuffling, and stochastic operations.

\subsection{Model Comparison}
\label{sec:comparison}

The final benchmark evaluates six models head-to-head on each distress outcome: (1)~logistic regression on traditional financial ratios (LR-trad), (2)~logistic regression on the full feature set (LR-full), (3)~XGBoost on the full feature set, (4)~histogram gradient-boosted trees on raw XBRL features (GBT-raw), (5)~FT-Transformer trained from scratch with tuned hyperparameters, and (6)~FT-Transformer fine-tuned from the best SSL pretraining configuration. All models are evaluated on the identical held-out test split using the metrics defined in \Cref{sec:metrics}.
